\section{Completion to full iPRA without a standard-model shortcut}
\label{sec:completion}

This section records the remaining internal arithmetic bridge, adapted
from the equality-reflection and primitive-recursion appendices in the
supplied corpus~\cite{generated}. Its input is the predecessor just proved
and the derived ordinary iterator $I=\widehat I$. Every operation below
is consequently a fixed finite SR term. The order of construction is
predecessor, actual equality reflection, quantifier-free induction,
pairing, and general primitive recursion. We do not assume a general
PR compiler in order to construct its prerequisites.

\subsection{A characteristic of actual equality}

Define
\[
\begin{gathered}
 C(n,a,b)=I(n,a;u.b),\qquad x+y=I(y,x;u.Su),\\
 d(x,y)=I(y,x;u.Pu),\qquad
 \operatorname{sg}(x)=C(x,0,1),\qquad W(x)=C(x,1,Px).
\end{gathered}
\]
They satisfy their ordinary zero/successor equations. Applying $P$ to
$Sx=Sy$ proves injectivity. Applying $Z$ to $Sx=0$ and using separation
proves that a successor is nonzero. Coverage now strengthens to
$x=0\lor x=S(Px)$. Its two branches prove
\[
\begin{gathered}
 \operatorname{sg}(x)\in\{0,1\},\qquad
 \operatorname{sg}(x)=0\ \longleftrightarrow\ x=0,\\
 W(x)=0\ \longleftrightarrow\ x=1,\qquad C(x,a,a)=a.
\end{gathered}
\]
The displayed finite-range notation abbreviates an ordinary disjunction;
it is a theorem obtained by cases, not an extra induction schema.

Four literal-equation inductions prove
\[
 0+y=y,\qquad (Sx)+y=S(x+y),\qquad
 d(Sx,Sy)=d(x,y),\qquad d(Sa,a)=1.
\]
The first two induct on $y$ using addition. For the third the base is
$P(Sx)=x$, and its step applies $P$ to the current induction equation.
For the fourth, induction on $a$ uses the third identity at its step.

With $x$ fixed, one further equation induction on $y$ proves
\[
 C\bigl(\operatorname{sg}(d(x,y)),x,y+d(x,y)\bigr)=x.
 \tag{Residual}\label{eq:completion-residual}
\]
At zero this is $C(\operatorname{sg}(x),x,x)=x$. At a successor put
$r=d(x,y)$ and $r'=Pr=d(x,Sy)$. If $r'=0$, the new selector returns
$x$. Otherwise coverage of $r$ and $P0=0$ force $r=Sr'$. The induction
equation gives $y+r=x$, while the addition equations give
$(Sy)+r'=y+Sr'=y+r=x$. This proves the new selected branch. The
invariant being inducted on is exactly the displayed single equation.

It follows that
\[
                    d(Sa,b)=1\quad\longleftrightarrow\quad a=b.
\]
For the forward direction, Residual gives $b+1=Sa$, hence $Sb=Sa$;
apply $P$. For the reverse direction substitute $b=a$ into $d(Sa,a)=1$.
Therefore
\[
 e(a,b)=\operatorname{sg}(W(d(Sa,b)))
\]
has values $0,1$ and satisfies $e(a,b)=0\leftrightarrow a=b$ for the
logic's actual equality throughout the source domain. Equality is now
decidable: the zero characteristic gives equality, and characteristic
one contradicts any alleged equality by separation.

\subsection{Every quantifier-free induction instance}

For each fixed quantifier-free formula define, by external recursion on
its finite syntax,
\[
\begin{aligned}
 b_{s=t}&=e(s,t),& b_\bot&=1,\\
 b_{A\land B}&=C(b_A,b_B,1),&
 b_{A\lor B}&=C(b_A,0,b_B),&
 b_{A\to B}&=C(b_A,b_B,0).
\end{aligned}
\]
Induction on formula syntax proves within the theory that $b_A$ has
values $0,1$ and $b_A=0\leftrightarrow A$. The conjunction and
disjunction cases are the two Boolean alternatives for $b_A$. For
implication, characteristic zero of $A$ supplies $A$, so the implication
has the characteristic of $B$. Characteristic one of $A$ proves
$\neg A$, so the implication is true by intuitionistic explosion and
its characteristic is zero. No excluded middle for quantified formulas
is used.

Given a contextual induction base and step for $A(n)$, convert them to
the base and step for the literal equation $b_A(n)=0$, use EI, and
convert back. Arbitrary retained quantified assumptions are permitted.
This obtains all quantifier-free induction in the scalar language.

\subsection{Pairing and the general recursion compiler}

The familiar semiring and linear-order identities now follow by
quantifier-free inductions from addition, predecessor, and monus. In
particular define $a\le b$ by $d(a,b)=0$ and
$a\cdot b=I(b,0;u.u+a)$. All these are scalar terms.

For fixed $n$ set
\[
 U_n(v)=C(d((Sv)^2,n),Sv,v),\qquad
 B_n(i)=I(i,0;v.U_n(v)),\qquad q(n)=B_n(n).
\]
Quantifier-free induction on $i$, keeping $n$ fixed, gives the simultaneous
invariant
\[
 B_n(i)\le i,\quad B_n(i)^2\le n,\quad
                  B_n(i)=i\ \lor\ n<(S B_n(i))^2.
\]
Its base is immediate. If the update increments, its test rules out
the strict-upper-bound alternative, so the old state is $i$ and the new
state is $Si$; the square bound is the test. If it does not increment,
the strict upper square bound holds and the old state is retained.
At $i=n$, the alternative $B_n(n)=n$ also implies the strict upper
square bound, since $n<(n+1)^2$. Consequently
\[
                       q(n)^2\le n<(q(n)+1)^2.
\]
Monotonicity of squares proves uniqueness of this root interval.

Define
\[
 \pi(a,b)=(a+b)^2+a,\quad
 \ell(n)=d(n,q(n)^2),\quad r(n)=d(q(n),\ell(n)).
\]
For $k=a+b$, $k^2\le k^2+a<(k+1)^2$ because $a\le k$.
Root uniqueness gives $q(\pi(a,b))=a+b$, and subtraction gives
\[
                 \ell(\pi(a,b))=a,\qquad r(\pi(a,b))=b.
\]
Surjectivity of $\pi$ is unnecessary; the domain is never restricted to
its range.

Suppose the earlier descriptions $g(\vec p)$ and $h(n,z,\vec p)$ have
already been compiled. Put
\[
 V(n,\vec p)=I\bigl(n,\pi(0,g(\vec p));
       u.\pi(S\ell(u),h(\ell(u),r(u),\vec p))\bigr),\qquad
 f(n,\vec p)=r(V(n,\vec p)).
\]
Equation induction proves $\ell(V(n,\vec p))=n$. The two projection
identities then give
\[
 f(0,\vec p)=g(\vec p),\qquad
 f(Sn,\vec p)=h(n,f(n,\vec p),\vec p).
\]
Initial functions and composition are immediate. External induction on
each fixed finite PR description therefore supplies its SR term and
internal defining-equation proofs. Every translated quantifier-free
induction instance is available by the preceding reflection argument.

\subsection{Both return interpretations}

Interpret SR in iPRA by primitive recursion with seed $Sa$, and interpret
each PR description in SR by the displayed compiler. In either composite,
first identify successor: $Sa=R(0,a;u.0)$ maps to ordinary $Sa$ by the
zero clause, and compiling that successor back gives the source $Sa$.
This anchors the successor operation before uniqueness is invoked.
Then,
compare a fixed operation with its double translation after its earlier
subdescriptions have been identified. Their bases and updates agree;
equation induction on the equality of their values proves their agreement
at every argument. Keep initializer and passive parameters independent
throughout that induction, and substitute dependent terms only afterward.
This applies equally to SR and ordinary PR recursion.

Both domains are the entire original numerical domain, defined by $x=x$;
both equality relations are actual equality; and each definable numerical
isomorphism has graph $x=y$. The ordinary connectives, falsity, and
quantifiers are preserved, so formula induction gives the full composites.
This establishes bi-interpretability, not merely mutual interpretation or
standard-model agreement. Effective finite proof translations preserve
contradiction in both directions and therefore establish equiconsistency.
